% ---------------------------------------------------------------------------
% dft_numbers.tex
%
% Every number in the manuscript that comes from the DFT+U campaign is defined
% here and nowhere else, in both the article and the Supplemental Material.
%
%   campaign (a)  results/dftu_row_ordered      row-ordered Li_{1/2}CoO2
%   campaign (b)  results/dftu_order_contrast   row / zigzag / clumped, x = 1/2
%   campaign (c)  results/dftu_x13              sqrt3 x sqrt3 Li_{1/3}CoO2
%   distillation  results/dftu_summary.json
%
% Usage convention: a macro is followed by "~" and then its unit, e.g.
% \comomentRow~$\mu_{\mathrm{B}}$.  "~" is in the xspace exception list.
%
% PROVENANCE NOTE (verification finding B7).  The x = 1/2 site moments, charge
% split and the moment ceiling derived from them are quoted from the ground
% state, afmlo_intra_U4.0 in cell B1, and not from the ferromagnetic solution of
% the seven-atom cell.  \comomentRowRange records the spread across the FM and
% both AFM configurations, which is the honest measure of how well determined
% the moment is.
%
% RETIRED MACROS, and why.  \deltaEOrder and the original reading of
% \dosContrast assumed the row motif lies below the zigzag motif and the ordered
% motif carries the smaller density of states at E_F; neither survived.
% \dosContrast now names the measured U = 0 ratio, which is the falsification
% itself.  \momentCellRow, \enmfmRow, \dosRow, \dosZigzag, \dosClumped,
% \dosXthird and \deltaMdft are retired, the last because the bulk ground states
% are antiferromagnetic and supply no net Delta M.  \jintraRow and \jinterRow
% are retired in favour of \eafmIntra and \eafmInter, which are signed energy
% differences and carry no implicit spin-Hamiltonian convention; \jeffXthird is
% retired in favour of \splitXthird, the raw splitting per cell.
% ---------------------------------------------------------------------------

\newcommand{\Uused}{\ensuremath{4.0}\xspace}

% --- (a) row-ordered Li_{1/2}CoO2, ground state AFM within a Co4+ row ------
\newcommand{\comomentRow}{\ensuremath{0.97}\xspace}       % mu_B on Co4+, ground state
\newcommand{\comomentThreeRow}{\ensuremath{0.00}\xspace}  % mu_B on Co3+, ground state
\newcommand{\comomentRowRange}{\ensuremath{0.97}\ to \ensuremath{1.02}\xspace} % across FM and both AFM
\newcommand{\msatRow}{\ensuremath{0.48}\xspace}           % mu_B per Co if every moment aligned
\newcommand{\dSbound}{\ensuremath{0.23}\xspace}           % k_B per Co, bulk reading of Eq. (S34)
\newcommand{\magorderRow}{antiferromagnetic within a Co\ensuremath{^{4+}} row\xspace}
\newcommand{\efmafmRow}{\ensuremath{28.2}\xspace}         % meV/Co, FM above AFM
\newcommand{\enmRow}{\ensuremath{63.7}\xspace}            % meV/Co, NM above AFM
\newcommand{\gapRow}{\ensuremath{0.70}\xspace}            % eV, AFM ground state (0.698)
\newcommand{\gapRowFM}{\ensuremath{0.43}\xspace}          % eV, FM in the same cell (0.428)
\newcommand{\eafmIntra}{\ensuremath{-56.5}\xspace}        % meV/Co4+, E(AFM) - E(FM), intra-row
\newcommand{\eafmInter}{\ensuremath{+41.7}\xspace}        % meV/Co4+, E(AFM) - E(FM), inter-row
\newcommand{\jintraMag}{\ensuremath{56}\xspace}           % meV, magnitude of the dominant coupling
\newcommand{\dnRow}{\ensuremath{0.20}\xspace}             % e, d-occupation split, ground state
\newcommand{\bondRow}{\ensuremath{0.001}\xspace}          % Angstrom, contrast of site-mean Co-O
\newcommand{\bondSpanRow}{\ensuremath{0.021}\xspace}      % Angstrom, span of individual Co-O bonds

% --- (b) three periodic Li motifs at x = 1/2, 4x2 cell, 8 f.u. ------------
%     Fixed ideal geometry, ferromagnetic alignment imposed.
\newcommand{\deltaEZigzag}{\ensuremath{-8.3}\xspace}       % meV/f.u. vs row, U=4
\newcommand{\deltaEZigzagUzero}{\ensuremath{+8.2}\xspace}  % meV/f.u. vs row, U=0
\newcommand{\deltaEClumped}{\ensuremath{58.1}\xspace}      % meV/f.u. vs row, U=4
\newcommand{\deltaEClumpedUzero}{\ensuremath{92.4}\xspace} % meV/f.u. vs row, U=0
\newcommand{\gapRowMotif}{\ensuremath{0.679}\xspace}       % eV
\newcommand{\gapZigzag}{\ensuremath{0.800}\xspace}         % eV
\newcommand{\gapClumped}{\ensuremath{0.675}\xspace}        % eV
\newcommand{\gapMotifSpread}{\ensuremath{0.13}\xspace}     % eV, 0.800 - 0.675
\newcommand{\dosRowUzero}{\ensuremath{22.5}\xspace}        % states/eV/cell, U=0
\newcommand{\dosContrast}{\ensuremath{1.004}\xspace}       % clumped/row at U=0
\newcommand{\momentRowUzeroA}{\ensuremath{0.3990}\xspace}  % mu_B, row cell at U=0, site 1
\newcommand{\momentRowUzeroB}{\ensuremath{0.3986}\xspace}  % mu_B, row cell at U=0, site 2

% --- (c) sqrt3 x sqrt3 Li_{1/3}CoO2, ground state ferrimagnetic -----------
\newcommand{\comomentXthird}{\ensuremath{1.04}\xspace}     % mu_B on each Co4+
\newcommand{\magorderXthird}{ferrimagnetic, the two Co\ensuremath{^{4+}} moments antiparallel\xspace}
\newcommand{\efmafmXthird}{\ensuremath{12.4}\xspace}       % meV/Co, FM above ferri
\newcommand{\enmXthird}{\ensuremath{230}\xspace}           % meV/Co, NM above ferri
\newcommand{\splitXthird}{\ensuremath{37}\xspace}          % meV per cell, FM above ferri
\newcommand{\gapXthird}{\ensuremath{1.21}\xspace}          % eV, ferri ground state (1.206)
\newcommand{\gapXthirdUfour}{\ensuremath{1.179}\xspace}    % eV, FM at U=4
\newcommand{\gapXthirdUfive}{\ensuremath{1.166}\xspace}    % eV, FM at U=5.5
\newcommand{\dosXthirdNM}{\ensuremath{7.7}\xspace}         % states/eV/cell, NM
\newcommand{\dnXthird}{\ensuremath{0.26}\xspace}           % e, d-occupation split
\newcommand{\bondXthird}{\ensuremath{0.003}\xspace}        % Angstrom, contrast of site-mean Co-O

% --- methodological numbers quoted in the Supplemental Material ------------
\newcommand{\comomentHighU}{\ensuremath{2.28}\xspace}      % mu_B, Co4+ at U=5.5, high spin
\newcommand{\momentFMa}{\ensuremath{1.028}\xspace}         % mu_B, inequivalent FM Co4+, supercell
\newcommand{\momentFMb}{\ensuremath{0.952}\xspace}         % mu_B, the other one
\newcommand{\forcePlateau}{\ensuremath{78}\xspace}         % meV/Angstrom, residual force plateau
\newcommand{\primSign}{\ensuremath{3.1}\xspace}            % meV/Co, primitive-cell FM above NM
